\documentclass[10pt,a4paper]{amsart}
\usepackage[margin=19mm]{geometry}
\usepackage{amsmath,amssymb,mathtools}
\usepackage[scr=rsfs,cal=euler]{mathalfa}
\usepackage{xfrac}
\usepackage[unicode,hidelinks,bookmarks=false]{hyperref}
\newcommand{\ie}{i.e.\ }
\newcommand{\cf}{cf.\ }
\newcommand{\resp}{respectively\ }
\newcommand{\Subsection}{\subsection}
\providecommand{\nobreakdash}{\nobreak\hbox{-}}
\setlength{\emergencystretch}{2em}
\multlinegap0pt
\usepackage{enumerate}
\usepackage{amssymb}
\usepackage{dsfont}
\newtheorem{theorem}{Theorem}
\newtheorem{lemma}{Lemma}
\title{Solvability of a class of evolution operators on compact Lie groups}
\author{Alexandre Kirilov, Wagner Augusto Almeida de Moraes, Pedro Meyer Tokoro}
\date{Source: arXiv:2602.15203v1 (CC BY 4.0)}
\begin{document}
\maketitle

\noindent\emph{Source and licence.} Solvability of a class of evolution operators on compact Lie groups. Version arXiv:2602.15203v1 (CC BY 4.0), distributed by arXiv under the Creative Commons Attribution 4.0 International licence, http://creativecommons.org/licenses/by/4.0/, as recorded in the arXiv licence metadata for this exact version and shown on the abstract page https://arxiv.org/abs/2602.15203v1. Full licence notice: \texttt{licenses/arxiv-2602.15203-CC-BY-4.0.txt}.\par\smallskip

\noindent\textit{Editorial scope.} Counted unit: the article's theorem thm:lambda0\_cases (source0.tex lines 941--964). The article's complete original proof of it is delivered with it (source0.tex lines 965--1065, one \texttt{proof} environment of 101 lines). No mathematics is written, repaired or re-derived: the statement and the proof are the article's own lines, in the article's own notation, and no retained line is substituted or re-flowed.  Retained uncounted as the source-local context the proof needs: lines 350--352; lines 367--390; lines 928--930; lines 1069--1081.  Nothing was omitted from any retained span, so the package carries no omission record.  No retained line contains a citation, so the wrapper carries no bibliography and no bibliography entry.  No retained line contains a figure environment, an image-inclusion command or drawing code, so no image is delivered and the package declares no asset dependency.  Editorial formatting only, in addition: an AMS \texttt{amsart} preamble that restates the article's own declarations and macros used by the retained lines, namely the environments \texttt{theorem}, \texttt{lemma} on separate counters, together with the standard packages those lines need.  The retained environments print in this document as Theorem 1, Lemma 1 in order of appearance, whereas the article prints the same items with the numbering of its own full document.  The complete native source of the article as distributed in the arXiv e-print for this exact version is bundled unchanged as \texttt{sources/194783/source0.tex}.
	\begin{equation}\label{eq:P_cv}
		Pu=Lu-(s(t)+i\delta q(t))u-\alpha q(t)\overline{u},
	\end{equation}

	\begin{theorem}\label{thm:main_varcoef}
		Suppose that the operator $P$ in \eqref{eq:P_cv} satisfies:
		\begin{itemize}
			\item[(I)] $|\alpha|\neq|\delta|$;
			\item[(II)] For all $[\xi]\in\widehat{G}$ and $1\le m\le d_\xi$, there is no integer $k\in\mathbb{Z}$ satisfying
			\[
			\begin{cases}
				\operatorname{Re}\bigl(A_0(2\pi k+\mu_m(\xi)\overline{C_0})\bigr)=0,\\[2mm]
				|2\pi k+\mu_m(\xi)C_0|^2=|A_0|^2-|B_0|^2;
			\end{cases}
			\]
			\item[(III)] There exists $M>0$ such that, for all $[\xi]\in\widehat{G}$ with $\langle\xi\rangle\ge M$,
			\[
			\min\left\{
			\left|e^{-\rho_m(\xi) q_0}-e^{i\mu_m(\xi)p_0 2\pi+s_0}\right|,
			\,
			\left|1-e^{-\rho_m(\xi) q_0+i\mu_m(\xi)p_0 2\pi+s_0}\right|
			\right\}\ge \langle\xi\rangle^{-M},
			\]
			for every $1\le m\le d_\xi$.
		\end{itemize}
		
		Then for every $f\in C^\infty(\mathbb{T}\times G)$ there exists $u\in C^\infty(\mathbb{T}\times G)$ such that $Pu=f$.
	\end{theorem}

	\begin{equation}\label{eq:P_lambda0}
		Pu=\partial_t u-p_0Xu-(s(t)+i\delta q(t))u-\alpha q(t)\overline{u},
	\end{equation}

	\begin{lemma}\label{lem:dc_equiv}
		The Diophantine condition \textup{(DC)} is equivalent to the following:
		\medskip\noindent
		
		\textup{(DC')} There exists $M > 0$ such that
		\[
		[\xi] \in \widehat{G},\ \langle\xi\rangle \ge M
		\ \Longrightarrow\
		\left|e^{i(\mu_m(\xi)p_0 2\pi - q_0\sqrt{\delta^2 - |\alpha|^2})} - 1\right|
		\ge \langle\xi\rangle^{-M},
		\qquad 1 \le m \le d_\xi.
		\]
	\end{lemma}

	\begin{theorem}\label{thm:lambda0_cases}
		Suppose that the operator $P$ in \eqref{eq:P_lambda0} satisfies one of the following conditions:
		\begin{itemize}
			\item[(1)] $|B_0| > |A_0|$;
			\item[(2)] $|B_0| \le |A_0|$, $|\alpha| > |\delta|$ and, for all $[\xi] \in \widehat{G}$ and $1 \le m \le d_\xi$, there is no integer $k \in \mathbb{Z}$ satisfying
			\begin{equation}\label{eq:lambda0_resonance}
				\begin{cases}
					\operatorname{Re}(A_0(k + \mu_m(\xi)p_0)) = 0, \\
					4\pi^2 |k + \mu_m(\xi)p_0|^2 = |A_0|^2 - |B_0|^2;
				\end{cases}
			\end{equation}
			\item[(3)] $|\alpha| < |\delta|$ and $s_0 \neq 0$;
			\item[(4)] $|\alpha| < |\delta|$, $s_0 = 0$, for all $[\xi] \in \widehat{G}$ and $1 \le m \le d_\xi$, there is no $k \in \mathbb{Z}$ satisfying \eqref{eq:lambda0_resonance}, and the following Diophantine condition holds:
			
			\medskip\noindent
			\textup{(DC)} There exists $M > 0$ such that for all $[\xi] \in \widehat{G}$ with $\langle\xi\rangle \ge M$,
			\[
			|2\pi k + \mu_m(\xi)p_0 2\pi - q_0\sqrt{\delta^2 - |\alpha|^2}| \ge \langle\xi\rangle^{-M},
			\qquad 1 \le m \le d_\xi.
			\]
		\end{itemize}
		
		Then, for every $f \in C^\infty(\mathbb{T} \times G)$, there exists $u \in C^\infty(\mathbb{T} \times G)$ such that $Pu = f$.
	\end{theorem}

	\begin{proof}
		Let $f \in C^\infty(\mathbb{T} \times G)$ and suppose that 
		$u \in \mathcal{D}'(\mathbb{T} \times G)$ satisfies $Pu=f$. 
		Proceeding as in the proof of Theorem~\ref{thm:main_varcoef}, 
		we take partial Fourier coefficients with respect to $G$ and 
		reduce the equation to a $2\times2$ system of ODEs for each 
		$[\xi] \in \widehat G$ and $1 \le m \le d_\xi$.
		
		When $\lambda=0$, the quantity
		\[
		\rho_m(\xi)=\sqrt{(-i\delta)^2+|\alpha|^2}
		=\sqrt{|\alpha|^2-\delta^2}
		\]
		is independent of $\xi$, and we simply write $\rho$.
		As in the general case, the solution is obtained by variation of constants,
		and periodicity leads to denominators of the form
		\[
		e^{-\rho q_0}-e^{i\mu_m(\xi)p_0 2\pi+s_0}
		\quad\text{and}\quad
		1-e^{-\rho q_0+i\mu_m(\xi)p_0 2\pi+s_0}.
		\]
		
		Thus, global solvability reduces to controlling these quantities.
		
		We first observe that if $|\alpha| \ne |\delta|$, then $\rho \ne 0$,
		so no degeneracy occurs in the diagonalization procedure.
		Moreover, the resonance system excluded in
		\eqref{eq:lambda0_resonance} corresponds precisely to the
		vanishing of the above denominators, as in the proof of
		Theorem~\ref{thm:main_varcoef}.
		
		If $|\alpha|>|\delta|$, then $\rho\in\mathbb R$.
		In this case $e^{-\rho q_0}$ is a positive real number.
		The identity
		\[
		e^{-\rho q_0}=e^{i\mu_m(\xi)p_0 2\pi+s_0+i2\pi k}
		\]
		would imply simultaneously that the imaginary part vanishes
		and that
		\[
		-\rho q_0 = s_0 + i2\pi(\mu_m(\xi)p_0+k),
		\]
		which yields exactly the system \eqref{eq:lambda0_resonance}.
		Hence, under the hypotheses excluding that system,
		the denominators do not vanish.
		Furthermore, if $|B_0|>|A_0|$, then
		$e^{-\rho q_0}\neq e^{s_0}$,
		so the distance between these exponentials is bounded below by
		a strictly positive constant independent of $[\xi]$.
		The same conclusion follows when $|B_0|\le |A_0|$ but
		\eqref{eq:lambda0_resonance} has no integer solution,
		since the trivial representation would otherwise produce a contradiction.
		In both situations we obtain a uniform bound
		\[
		\left|e^{-\rho q_0}-e^{i\mu_m(\xi)p_0 2\pi+s_0}\right|
		+
		\left|1-e^{-\rho q_0+i\mu_m(\xi)p_0 2\pi+s_0}\right|
		\ge C>0,
		\]
		independent of $[\xi]$.
		
		If instead $|\alpha|<|\delta|$, then $\rho=i\omega$ with
		$\omega=\sqrt{\delta^2-|\alpha|^2}\in\mathbb R$,
		so $|e^{-\rho q_0}|=1$.
		If $s_0\neq0$, then
		\[
		|e^{-\rho q_0}-e^{i\mu_m(\xi)p_0 2\pi+s_0}|
		\ge |1-e^{s_0}|,
		\]
		which is strictly positive.
		Hence the denominators are again uniformly bounded away from zero.
		
		The remaining situation occurs when $|\alpha|<|\delta|$ and $s_0=0$.
		In this case the denominators reduce to
		\[
		|e^{i(\mu_m(\xi)p_0 2\pi-q_0\omega)}-1|,
		\]
		up to harmless sign changes.
		They may approach zero only if the phase
		$\mu_m(\xi)p_0 2\pi-q_0\omega$
		approaches an integer multiple of $2\pi$.
		The Diophantine condition (DC) prevents this phenomenon
		by imposing the lower bound
		\[
		|2\pi k+\mu_m(\xi)p_0 2\pi-q_0\omega|
		\ge \langle\xi\rangle^{-M},
		\]
		which is equivalent, by Lemma~\ref{lem:dc_equiv},
		to a polynomial lower bound on
		$|e^{i(\mu_m(\xi)p_0 2\pi-q_0\omega)}-1|$.
		Consequently, the inverses of the denominators grow at most
		polynomially in $\langle\xi\rangle$.
		
		In all cases, therefore, the constants arising from periodicity
		either remain uniformly bounded or have at most polynomial growth.
		Since the forcing terms are rapidly decreasing in $\xi$,
		the same estimates as in Theorem~\ref{thm:main_varcoef}
		show that the Fourier coefficients of $u$
		decay rapidly.
		Hence $u\in C^\infty(\mathbb{T}\times G)$ and $Pu=f$.
	\end{proof}

\end{document}
